% ===================================================================
%  TABELL Nr. 6 — D'Haus bannen, d'Miwwelen, d'Iessen, d'Beliichtung.
%  Traduction 2026 d'apres texto/io, controlee sur texto/fr, texto/de
%  et texto/nl. Consigne : voir texto/lb/10-tabell-01.tex.
%
%  LE PLUS LONG TABLEAU DU LIVRET, cinq pieces et plus de deux cents
%  objets, et il verifie piece par piece ce que le tableau 5 avait
%  trouve sur le batiment : LA LIGNE ENTRE LES DEUX VOISINS PASSE A
%  L'INTERIEUR DE LA MAISON.
%
%  La cuisine ne coute pas un mot francais : « Häerd » l'atre,
%  « Blosbalg » le soufflet, « Kachuewen » le fourneau,
%  « Hackbriet » le billot, « Biesem » le balai, « Spülstee »
%  l'evier, « Wasserhunn » le robinet, « Krou » la cruche,
%  « Bidden » le baquet, « Sief » le tamis, « Riffelen » la rape.
%
%  Le salon n'en coute presque que : « Lambrequin », « Sekretär »,
%  « Rideauen », « Embrassen », « Paravent », « Portrait »,
%  « Album », « Kanapee », « Fauteuil », « Luster », « Sonnerie »,
%  « Tabouret », « Piano ».
%
%  ET LA SALLE A MANGER EST EXACTEMENT AU MILIEU, ce qui n'etait pas
%  prevu. On y mange avec des mots germaniques — « Teller »,
%  « Läffel », « Messer », « Glas », « Brout », « Kéis » — et l'on y
%  recoit avec des mots francais — « Menagère », « Portière »,
%  « Buffet », « Serviette », « Karaff », « Dessert ». La table est
%  allemande, le service est francais.
%
%  C'EST DONC BIEN L'USAGE ET NON LA PIECE QUI DECIDE. Le tableau 5
%  avait laisse croire que la ligne separait des salles ; elle
%  separe, dans la meme salle, LE GESTE DE LA MAIN DU GESTE DE LA
%  POLITESSE. On mange dans la langue de la cuisine et l'on sert dans
%  celle du salon.
% ===================================================================

%%K t06-apar-1 apar
\VUsaut{6.0mm}
\VUtitre{15.0pt}{60}{TABELL Nr. 6}
\VUsaut{1.4mm}
\VUtitre{11.4pt}{0}{D'Haus bannen, d'Miwwelen, d'Iessen,}
\VUsaut{0.4mm}
\VUtitre{11.4pt}{0}{d'Beliichtung.}
\VUsaut{2.2mm}

%%K t06-apar-2 apar
D'Haus ass fäerdeg. Dem Jang säi Monni ass a senger neier Wunneng agezunn, vun där
mir d'Antelung kennen. Elo kucke mer, wéi hie säin Haus ameubléiert huet. Fir
d'éischt entdecke mer:

%%K t06-tit-1 sub
\VUpk{10.2pt}{40}{D'Schlofzëmmer.}
\VUsaut{3.0mm}

%%K t06-c1-01-1 p
1. --- Mir kommen zur Onzäit, well d'Zëmmermeedchen ass amgaangen d'\VUgras{Zëmmer}
ze botzen.

%%K t06-c1-02-1 p
2. --- Um \VUgras{Toilettendësch} \textsuperscript{(1)} leien hei an do verschidde
Saachen, mat deenen ee sech wäscht, kämmt, kuerz gesot sech fäerdeg mécht. Dat sinn
d'\VUgras{Kuvette} \textsuperscript{(2)} an d'\VUgras{Waasserkrou}
\textsuperscript{(3)}, d'\VUgras{Kappbiischt} \textsuperscript{(4)}, de
\VUgras{Kamm} \textsuperscript{(5)}, d'\VUgras{Seef} \textsuperscript{(6)} an de
\VUgras{Schwamm} \textsuperscript{(7)}, a schlussendlech e \VUgras{Fläschchen}
\textsuperscript{(8)} fir dat flëssegt Zännwaasser.

%%K t06-c1-03-1 p
3. --- Um Buedem, virum Toilettendësch, do ass den \VUgras{Eemer}
\textsuperscript{(9)}, fir dat dreckegt Waasser opzefänken.

%%K t06-c1-04-1 p
4. --- Am \VUgras{Virzëmmer} \textsuperscript{(10)} gesi mer e puer
\VUgras{Kleederhaken} \textsuperscript{(13)} an zwee \VUgras{Prabbelien}
\textsuperscript{(11)} am \VUgras{Prabbelisstänner} \textsuperscript{(12)}.

%%K t06-c1-05-1 p
5. --- Op der anerer Säit vun der Dier ass de \VUgras{Spigelschaf}
\textsuperscript{(14)}, an deem d'\VUgras{Wäschzäich} \textsuperscript{(15)} ass.

%%K t06-c1-06-1 p
6. --- Notze mer aus, datt d'\VUgras{Zëmmermeedchen} \textsuperscript{(6)}
d'\VUgras{Bett} \textsuperscript{(17)} mécht, fir seng verschidden Deeler ze
entdecken. D'\VUgras{Bettgestell} \textsuperscript{(20)} hält eng gutt
\VUgras{Fiedermatratz} \textsuperscript{(21)}, op déi d'Justine gläich eng wullen
\VUgras{Matratz} \textsuperscript{(22)} leet. Duerno hëlt si vun de Still déi zwee
\VUgras{Bettdicher} \textsuperscript{(23)} an d'\VUgras{Decken}
\textsuperscript{(24)}. Da leet si um Kappenn dat \VUgras{Nackekëssen}
\textsuperscript{(25)} an dat \VUgras{Kappkëssen} \textsuperscript{(26)}, déi mat
dem \VUgras{Duvet} \textsuperscript{(27)} op der \VUgras{Bettvirlag}
\textsuperscript{(32)} leien. Si botzt mat engem Fiederbësem d'\VUgras{Rideauen}
\textsuperscript{(19)} an den \VUgras{Betthimmel} \textsuperscript{(18)}, an do ass en
Deel vun der Aarbecht fäerdeg.

%%K t06-c1-07-1 p
7. --- Da muss si de \VUgras{Nuetsdësch} \textsuperscript{(28)} e bësse riicht
maachen, de \VUgras{Wecker} \textsuperscript{(29)} opzéien, dat
\VUgras{Nuetslämpchen} \textsuperscript{(30)} virbereeden an d'\VUgras{Kerz}
\textsuperscript{(31)} erausthuelen, fir de Kerzestänner ze botzen.

%%K t06-tit-2 sub
\VUsaut{5.0mm}
\VUpk{10.2pt}{40}{De Bidden a d'Kamäinsgeschir.}
\VUsaut{3.0mm}

%%K t06-c1-08-1 p
8. --- De \VUgras{Bidden} \textsuperscript{(34)} nieft dem \VUgras{Tabouret}
\textsuperscript{(33)} brauch och ganz noutwenneg opgeraumt ze ginn. Si muss
d'\VUgras{Stickereien} \textsuperscript{(35)} zesummeleeën, an den
\VUgras{Aarbechtsdësch} \textsuperscript{(38)} de \VUgras{Fangerhutt,} de
\VUgras{Fuedem} \textsuperscript{(36)}, d'\VUgras{Nolen} \textsuperscript{(37)} an
d'\VUgras{Schéier} \textsuperscript{(39)} eraleeën, an de Deckel vum Dësch mat dem
\VUgras{Schlëssel} \textsuperscript{(40)} zoumaachen.

%%K t06-c1-09-1 p
9. --- Um \VUgras{Guéridon} \textsuperscript{(41)} an der Mëtt wäert d'Justine
d'Waasser vun der \VUgras{Karaff} \textsuperscript{(42)} erneieren. Si wäert
\VUgras{Zocker} an d'\VUgras{Zockerdous} \textsuperscript{(43)} leeën, d'\VUgras{Zang}
\textsuperscript{(44)} botzen an d'\VUgras{Kleederbiischt} \textsuperscript{(46)}
erofhuelen.

%%K t06-c1-10-1 p
10. --- Déi riets Säit vum Zëmmer wäert manner Aarbecht vun hir verlaangen, well de
\VUgras{laange Fauteuil} \textsuperscript{(47)} a säi \VUgras{Këssen}
\textsuperscript{(48)} sinn op hirer richteger Plaz, a well et net kal ass, ass näischt
tëscht dem Geschir vum \VUgras{Kamäin} \textsuperscript{(49)} verréckelt.
D'\VUgras{Schëpp} \textsuperscript{(50)}, d'\VUgras{Zaang} \textsuperscript{(51)},
d'\VUgras{Feiereisen} \textsuperscript{(55)} an d'\VUgras{Äschegitter}
\textsuperscript{(56)} si propper; de \VUgras{Kamäinsschirm} \textsuperscript{(54)}
gouf net verréckelt, an d'\VUgras{Holzkëscht} \textsuperscript{(52)} ass gutt mat
\VUgras{Holz} \textsuperscript{(53)} versuergt.

%%K t06-noto-1 noto
\VUnotes{143.2mm}{%
(*) Babo = klengt Kand; engl. \textit{baby}, franz. \textit{bébé}, ital.
\textit{bambino}.%
}

%%K t06-noto-2 noto
\VUnotes{143.2mm}{%
(*) Lambrequino = franz. \textit{lambrequin}, span. \textit{lambrequines}, port.
\textit{lambrequins}, souguer däitsch an engl. \textit{lambrequins}, also däitsch
engl. franz. port. span.%
}

%%K t06-c1-11-1 p
11. --- Trotzdeem muss si d'\VUgras{Pendul} \textsuperscript{(57)}, d'\VUgras{Girandolen}
\textsuperscript{(58)} an d'\VUgras{Dousen} \textsuperscript{(59)} ofstëpsen, a
schlussendlech de \VUgras{Spigel} \textsuperscript{(60)} ofwëschen.

%%K t06-c1-12-1 p
12. --- Wat d'\VUgras{Wéi} \textsuperscript{(61)} vum klenge Kand (vum Babo (*))
ubelaangt, déi ass scho prett.

%%K t06-tit-3 sub
\VUsaut{5.0mm}
\VUpk{10.2pt}{40}{De Salon.}
\VUsaut{3.0mm}

%%K t06-c2-01-1 p
1. --- Elo, do ass de Salon. Dat ass e ganz schéint \VUgras{Zëmmer.} De Kamäin, dee
riets am Eck steet, ass mat enger feiner \VUgras{Statuett} \textsuperscript{(1)} a mat
zwou schéine \VUgras{Vasen} \textsuperscript{(3)} geschmückt a mat engem sammetenen
\VUgras{Lambrequin} \textsuperscript{(2)} (*) drapéiert.

%%K t06-c2-02-1 p
2. --- Fir ze verhënneren, datt d'Fonken aus dem Häerd den Teppech ubrennen, huet
ee virum Häerd e koperen \VUgras{Fonkeschirm} \textsuperscript{(4)} gestallt.

%%K t06-c2-03-1 p
3. --- Nieft drun ass en elegante \VUgras{Sekretär} \textsuperscript{(5)} fir Dammen
op. Do drop schreift d'\VUgras{Hausdamm} \textsuperscript{(23)} hir Bréiwer.

%%K t06-c2-04-1 p
4. --- D'Fënster ass mat \VUgras{Fënsterrideauen} \textsuperscript{(6)} mat
\VUgras{Embrassen} \textsuperscript{(8)} un de \VUgras{Patèren} \textsuperscript{(7)}
geschmückt a mat Stoffrideaue mat \VUgras{Draperie uewen}
\textsuperscript{(9)}.

%%K t06-c2-05-1 p
5. --- An der Fënsternisch ass an engem \VUgras{Blummendëppen}
\textsuperscript{(11)} eng gréng \VUgras{Planz} \textsuperscript{(10)}, eng herrlech
Palm, vun där d'Blieder sech bal bis bei d'\VUgras{Pëtrolslamp}
\textsuperscript{(12)} strecken.

%%K t06-c2-06-1 p
6. --- De \VUgras{Lampeschirm} \textsuperscript{(13)} vun der Lamp gouf vun der
Maria geriicht, dem charmante jonke Meedchen \textsuperscript{(14)}, dat um Piano
\textsuperscript{(15)} sëtzt. Ganz riicht op dem \VUgras{Tabouret}
\textsuperscript{(16)} sëtzend, studéiert si e Museksstéck. Si ass ganz geschéckt a
suergfälteg, wéi hire \VUgras{Notenschaf} \textsuperscript{(17)} beweist, dee perfekt
opgeraumt ass.

%%K t06-tit-4 sub
\VUsaut{5.0mm}
\VUpk{10.2pt}{40}{De Bengel.}
\VUsaut{3.0mm}

%%K t06-c2-07-1 p
7. --- Hire Brudder, de Paul, sicht e \VUgras{Buch} \textsuperscript{(19)} an der
\VUgras{Bibliothéik} \textsuperscript{(18)}. Hien ass e \VUgras{Jonken}
\textsuperscript{(20)}, onroueg an onerdräiglech. Well hie sech un d'Bibliothéik
hänkt, fir dat Buch ze kréien, dat ze héich ass, riskéiert hien, de \VUgras{Bust}
\textsuperscript{(21)} erofzegeheien, deen uewe steet.

%%K t06-c2-08-1 p
8. --- Fir dës Dommheet ze maachen, notzt hien aus, datt seng Mamm amgaange schwätzt
mat enger Frëndin, enger \VUgras{Damm} \textsuperscript{(22)}, déi e puer Deeg bei hir
verbréngt. D'Mamm sëtzt op engem \VUgras{Kanapee} \textsuperscript{(24)} nieft dem
\VUgras{Fauteuil} \textsuperscript{(25)}, an hir Frëndin op engem \VUgras{Stull}
\textsuperscript{(27)}, vun deem ee nëmmen d'\VUgras{Lehn} \textsuperscript{(26)}
gesäit.

%%K t06-c2-09-1 p
9. --- De grousse \VUgras{Kader} \textsuperscript{(30)}, deen un der Mauer hänkt,
huet de \VUgras{Portrait} \textsuperscript{(29)} vun enger Verwandter. Am
\VUgras{Album} \textsuperscript{(32)}, deen um Dësch läit, sinn d'\VUgras{Fotoen}
\textsuperscript{(33)} vun den anere Familljemembere gesammelt. D'Kanner hunn en grad
duerchgeblieddert, well d'\VUgras{Still} \textsuperscript{(31)} stinn nach ronderëm de
Dësch zesummen.

%%K t06-c2-10-1 p
10. --- Déi porzelanen \VUgras{chinesesch Vas} \textsuperscript{(34)}, déi nieft dem
\VUgras{Pabeiersmesser} \textsuperscript{(35)} steet, gouf der Maria fir hire
Gebuertsdag geschenkt, an de \VUgras{Paravent} \textsuperscript{(36)}, deen den Eck
vum Zëmmer schmückt, gouf vun hirem Monni aus China matbruecht.

%%K t06-c2-11-1 p
11. --- Vun de wonnerschéinen \VUgras{Biller} \textsuperscript{(37)} opgehellt, déi un
der \VUgras{Mauer} \textsuperscript{(42)} hänken, vum \VUgras{Plafongsluster}
\textsuperscript{(38)} beliicht, vun deem d'\VUgras{Elektreschlampen}
\textsuperscript{(39)} owes frou blénken, gesäit dëse Salon ganz agreabel aus. De
mëllen \VUgras{Teppech} \textsuperscript{(40)}, dee sech iwwert de Parquet spreet, an
déi esou bequem \VUgras{elektresch Sonnerie} \textsuperscript{(41)} maachen säi
Komfort komplett.

%%K t06-tit-5 sub
\VUsaut{3.6mm}
\VUpk{10.2pt}{40}{D'Iesszëmmer.}
\VUsaut{3.0mm}

%%K t06-c3-01-1 p
1. --- D'Iessensklack huet gelaut. Déi ganz Famill, ausser der Maria, ass ronderëm
den Dësch versammelt. D'Iesszëmmer gëtt agreabel vum ganz gudde faience-enen
\VUgras{Uewen} \textsuperscript{(1)} mat sengem dënne \VUgras{Rouer}
\textsuperscript{(2)} gewiermt.

%%K t06-c3-02-1 p
2. --- Dëst Zëmmer huet net vill Miwwelen. Laanscht d'Mauer hänkt, iwwert der
\VUgras{Konsol} \textsuperscript{(4)}, e schéint \VUgras{Fontänchen}
\textsuperscript{(3)}. Ganz no drun ass de \VUgras{Buffet} \textsuperscript{(5)}, op
deen ee kalt Iessen, d'Dessertteller an déi verschidde Geschir stellt, déi een net
direkt brauch.

%%K t06-c3-03-1 p
3. --- Esou gesi mer um ieweschte Bord e \VUgras{gebrode Poulet}
\textsuperscript{(7)}, e \VUgras{Fësch} \textsuperscript{(8)} an eng \VUgras{Schuel}
\textsuperscript{(10)} mat \VUgras{Uebst} \textsuperscript{(9)}. Drënner, do sinn de
\VUgras{Kéis} \textsuperscript{(11)}, d'\VUgras{Brout} \textsuperscript{(12)}, d'
\VUgras{Menagère} \textsuperscript{(13)}, an där alles ass, wat ee brauch, fir de
Zalot unzemaachen: d'Ueleg, den Esseg, d'\VUgras{Salz} \textsuperscript{(14)} an de
\VUgras{Peffer} \textsuperscript{(15)}. Schlussendlech ass um ënneschte Bord e
\VUgras{Kuch} \textsuperscript{(17)} nieft dem \VUgras{Moschterdëppchen}
\textsuperscript{(16)}.

%%K t06-c3-04-1 p
4. --- D'Auer vum Iesszëmmer, déi ee \VUgras{Wanduhr} \textsuperscript{(20)} nennt,
huet eng ganz besonnesch Form. Si ass an en hëlzene Kader agebaut, deen ausserdeem e
\VUgras{Barometer} \textsuperscript{(19)} an e \VUgras{Thermometer}
\textsuperscript{(18)} huet.

%%K t06-tit-6 sub
\VUsaut{4.6mm}
\VUpk{10.2pt}{40}{De Service beim Iessen.}
\VUsaut{3.0mm}

%%K t06-c3-05-1 p
5. --- Un all Maueren am Zëmmer ass eng \VUgras{Holzverkleedung}
\textsuperscript{(21)} aus gewuessenem Eechenholz. Eng stoffend \VUgras{Portière}
\textsuperscript{(22)} mécht d'Dier zou, déi op d'Veranda geet. Do wäert d'Famill nom
Iesse Kaffi drénke goen. Schonn sinn d'\VUgras{Taassen} \textsuperscript{(24)} an
d'\VUgras{Ënnertaassen} \textsuperscript{(25)} op der \VUgras{Anrichte}
\textsuperscript{(23)} prett.

%%K t06-c3-06-1 p
6. --- Iwwert dem Dësch ass um Plafong eng \VUgras{Hängelamp}
\textsuperscript{(26)} befestegt, vun där déi ganz hell Luucht owes dat ganzt
Iesszëmmer beliicht.

%%K t06-c3-07-1 p
7. --- Elo entdecke mer den \VUgras{Iessdësch} \textsuperscript{(27)} selwer. Wéi
d'Famill erakoum, war d'Gedeck prett. Um \VUgras{Dëschduch} \textsuperscript{(28)}
waren, op der Plaz vun all Dëschgenoss, en \VUgras{Teller} \textsuperscript{(33)},
eng \VUgras{Serviette} \textsuperscript{(29)}, e \VUgras{Läffel}
\textsuperscript{(34)}, eng \VUgras{Forschett} \textsuperscript{(35)}, e
\VUgras{Messer} \textsuperscript{(36)} an e \VUgras{Glas} \textsuperscript{(32)}
geluecht. Et waren och \VUgras{Fläschen} \textsuperscript{(30)} fir de Wäin an eng
\VUgras{Karaff} \textsuperscript{(31)} fir d'Waasser.

%%K t06-c3-08-1 p
8. --- De \VUgras{Bedéngten} \textsuperscript{(39)} huet fir d'éischt d'\VUgras{Zoppeschossel}
\textsuperscript{(37)} bruecht, an där nach e bësse Zopp dämpft. Elo servéiert hien
dee éischte \VUgras{Plat} \textsuperscript{(38)}, e Fleeschplat. Duerno wäert hien déi
servéieren, déi mir scho genannt hunn, a schlussendlech, nom \VUgras{Dessertplat,}
wäert hien ausnotzen, datt seng Häre schonn op der Veranda sinn, fir d'Dëschgeschir
ewechzehuelen an d'Iesszëmmer nees opzeraumen.

%%K t06-c3-08-2 p
\textit{Bemierkung.} --- \textit{Déi geleefeg Wierder, déi d'Iessen ubelaangen, ginn
hei ganz kuerz uginn. D'Lëscht gëtt op deenen zwou Tabellen „De Viischtbutteck“
(Nr.~13) an „De Maart“ (Nr.~15) komplettéiert.}

%%K t06-tit-7 sub
\VUsaut{4.6mm}
\VUpk{10.2pt}{40}{D'Kichen.}
\VUsaut{3.0mm}

%%K t06-c4-01-1 p
1. --- An der Kichen, an déi mer duerch déi hallef op Dier kucken, gesi mer e
molerescht Duercherneen. Virum \VUgras{Häerd} \textsuperscript{(1)}, an deem e
\VUgras{Feier} \textsuperscript{(3)} knascht, ass e \VUgras{Bratspiessdréier}
\textsuperscript{(5)} installéiert. E schéine \VUgras{Broot} \textsuperscript{(7)} ass
\VUgras{opgespiisst} \textsuperscript{(6)} a brëtselt.

%%K t06-c4-02-1 p
2. --- De \VUgras{Blosbalg} \textsuperscript{(8)} an d'\VUgras{Kuelenschëpp}
\textsuperscript{(9)}, déi ee grad benotzt huet, sinn um Buedem bliwwen, sou wéi
d'\VUgras{Holz} \textsuperscript{(10)}.

%%K t06-c4-03-1 p
3. --- Uewen um Kamäin ass alles opgeraumt: d'\VUgras{Luucht}
\textsuperscript{(11)}, d'\VUgras{Zünnhëlzerdous} \textsuperscript{(13)}, an där
d'\VUgras{Zünnhëlzer} \textsuperscript{(12)} sinn, de \VUgras{Kerzestänner}
\textsuperscript{(14)} an de \VUgras{Kerzeliichter} \textsuperscript{(15)} mat der
\VUgras{Kerz} \textsuperscript{(16)} sinn op hirer Plaz. Mä de grousse
\VUgras{Kuelenemer} \textsuperscript{(18)}, voll \VUgras{Kuelen}
\textsuperscript{(17)}, déi d'\VUgras{Kächin} \textsuperscript{(23)} grad am Keller
gefëllt huet, ass an der Mëtt vun der Kichen bliwwen.

%%K t06-c4-04-1 p
4. --- Wierklech muss déi aarm Fra sech beeilen, fir datt d'Iesse rechtzäiteg prett
ass. Elo ass si beim \VUgras{Kachuewen} \textsuperscript{(19)} a réiert eng
\VUgras{Zooss} \textsuperscript{(24)}, déi ee net ze vill ubrenne loosse soll.

%%K t06-c4-05-1 p
5. --- Am \VUgras{Bakuewen} \textsuperscript{(20)} baacht e Kuch, dee si fir
d'Kanner hir Vesper virbereet huet. Iwwert dem Kachuewen hänken de
\VUgras{Schäffchen} \textsuperscript{(21)} an de \VUgras{Schaumläffel}
\textsuperscript{(22)}.

%%K t06-tit-8 sub
\VUsaut{4.0mm}
\VUpk{10.2pt}{40}{D'Geschir.}
\VUsaut{3.0mm}

%%K t06-c4-06-1 p
6. --- Dat \VUgras{Kachgeschir} \textsuperscript{(25)}, dat hannen am Zëmmer
opgehaange läit, ass ganz propper. Déi schéi koper \VUgras{Kasserollen}
\textsuperscript{(26)}, d'\VUgras{Mëllechdëppen} \textsuperscript{(27)}, de
\VUgras{Sief} \textsuperscript{(28)} an d'\VUgras{Riffelen} \textsuperscript{(29)}
goufe suergfälteg gebotzt.

%%K t06-c4-07-1 p
7. --- D'\VUgras{Salzkëscht} \textsuperscript{(30)} steet um Geschirschaf nieft der
\VUgras{Téikrou} \textsuperscript{(31)} an der \VUgras{Uelegfläsch}
\textsuperscript{(32)}. An der \VUgras{Schuff} \textsuperscript{(33)} sinn
wahrscheinlech d'Bestecker an d'Kichemesseren.

%%K t06-c4-08-1 p
8. --- De \VUgras{Biesem} \textsuperscript{(34)} an de \VUgras{Fiederbësem}
\textsuperscript{(35)} stinn an engem Eck. D'Kächin huet grad de \VUgras{Holzblock}
\textsuperscript{(36)} benotzt, si huet en nieft der Fënster gelooss.

%%K t06-c4-09-1 p
9. --- En \VUgras{Handduch} \textsuperscript{(37)} hänkt un der Mauer, nieft dem
\VUgras{Trichter} \textsuperscript{(38)}. An deem Deel vun der Kichen sinn de
\VUgras{Grill} \textsuperscript{(39)}, dat \VUgras{Hackmesser} \textsuperscript{(40)}
an dat \VUgras{Hackbriet} \textsuperscript{(41)} ewechgestallt. Duerno, iwwert dem
Hackmesser, op engem \VUgras{Bord} \textsuperscript{(42)}, sinn \VUgras{Dëppen}
\textsuperscript{(43)}, dat \VUgras{Kachdëppen} \textsuperscript{(44)} mat sengem
\VUgras{Deckel} \textsuperscript{(45)} an d'\VUgras{Marmit} \textsuperscript{(46)}.
D'\VUgras{Panne} \textsuperscript{(47)} hänkt nieft drun.

%%K t06-c4-10-1 p
10. --- Nëmmen dee koperen \VUgras{Wasserhunn} \textsuperscript{(48)} blénkt an
dësem Eck. De \VUgras{Spülstee} \textsuperscript{(49)}, op deem déi déck
\VUgras{Krou} \textsuperscript{(50)} steet, ass wahrscheinlech ganz bequem mat sengem
klenge Schaf, an deem ee dat \VUgras{Essegfässchen} \textsuperscript{(51)} mat sengem
\VUgras{Hunn} \textsuperscript{(52)}, d'\VUgras{Offallkëscht} \textsuperscript{(53)}
an déi \VUgras{dreckeg Telleren} \textsuperscript{(54)} hält.

%%K t06-tit-9 sub
\VUsaut{4.0mm}
\VUpk{10.2pt}{40}{Aner Saachen, déi an der Kichen stinn.}
\VUsaut{3.0mm}

%%K t06-c4-11-1 p
11. --- De \VUgras{Bidden} \textsuperscript{(55)}, an deem ee d'\VUgras{Geméis}
\textsuperscript{(58)} wäscht, an de gusseisene \VUgras{Kessel}
\textsuperscript{(56)}, dee um \VUgras{Fliesebuedem} \textsuperscript{(57)} steet, hu
wahrscheinlech och hir Plaz an deem Schaf.

%%K t06-c4-12-1 p
12. --- Der Kächin feelen sécher keng Uewen, well do ass nach, um Eck vum Dësch, e
\VUgras{Gasuewen} \textsuperscript{(59)}, op deem an enger klenger Kasserol Waasser
waarm gëtt. De \VUgras{Damp} \textsuperscript{(60)}, deen erauskënnt, weist eis, datt
et wahrscheinlech kacht. Wahrscheinlech gëtt dëst Waasser fir de Kaffi benotzt, well
do sinn, um anere Enn vum Dësch, d'\VUgras{Kaffiskrou} \textsuperscript{(64)} an
d'\VUgras{Kaffismillen} \textsuperscript{(63)}.

%%K t06-c4-13-1 p
13. --- Den Uewen ass mat dem \VUgras{Gasbrenner} \textsuperscript{(62)} duerch e
laangen \VUgras{Gummischlauch} \textsuperscript{(61)} verbonnen.

%%K t06-c4-14-1 p
14. --- D'\VUgras{Deegléinerin} \textsuperscript{(66)}, déi der Kächin hëllefe
kënnt, bereet d'Geméis fir d'Iessen vir. Wann et gebotzt a gewäsch ass, leet si et an
de \VUgras{Baseng} \textsuperscript{(65)} an d'waart op d'Stonn fir et ze kachen.

%%K t06-tit-10 sub
\VUsaut{4.0mm}
{\centering\fontsize{10.2pt}{10.2pt}\selectfont\textls[40]{\scshape D'Buedzëmmer.} \textit{(Buedstuff.)}\par}
\VUsaut{3.0mm}

%%K t06-c5-01-1 p
1. --- Eis bleift nach eng Onbescheidenheet ze maachen, fir déi wichtegst Stuffe vun
der Wunneng kennen ze léieren: d'Anlag vum Buedzëmmer ze gesinn. Dat brauch net vill
Zäit, well et ass net grouss, a mir wësse séier, datt d'\VUgras{Buedwan}
\textsuperscript{(1)} eng gutt \VUgras{Buedheizung} \textsuperscript{(2)} huet, datt
d'\VUgras{Seefeschuel} \textsuperscript{(3)} vun der Hand vum Bueder z'erreechen ass,
an datt den \VUgras{Handduchhalter} \textsuperscript{(5)} et sécher liicht mécht,
d'\VUgras{Handdicher} \textsuperscript{(4)} ze dréchnen.

%%K t06-c5-02-1 p
2. --- Dëst Zëmmer huet seng Maueren mat \VUgras{Fliesen} \textsuperscript{(6)}
bedeckt, déi et méiglech gemaach hunn, en \VUgras{Duschapparat}
\textsuperscript{(7)} anzeriichten, ouni Angscht ze hunn, datt d'Waasser, dat beim
Benotze spritzt, d'Maueren beschiedegt. E zinkent \VUgras{Beckelchen}
\textsuperscript{(8)}, dat fir d'Foussbieder déngt, mécht d'Miwwelen komplett.

\VUsaut{6.0mm}
\VUfilet{22mm}
